Spacing collapse, divergence of close-pair counts, and this empty local limit remain true under any increasing degree-one unfolding \(F_N\) satisfying \(\|F_N'\|_\infty\leq C N^A\) for fixed \(C,A\). Thus the present conclusion cannot be removed by an unfolding with polynomially bounded distortion. No assertion is made about unfoldings whose derivatives grow faster than every power of \(N\).

\subsection{Quantization and a finite locality estimate}\label{locality:quantization-and-a-finite-locality-estimate}

For an integer frequency \(v=(a,b)\) use the \textbf{physical integer Fourier lift} \[
 (Q_N(a,b)u)(x)=e^{2\pi ia(x-b/2)/N}u(x-b\bmod N).
\] The half in this exponent is ordinary division, not the inverse of two in \(\mathbb Z/N\mathbb Z\). For the modular-centred convention \(W_N\) used in the earlier notes, \(Q_N(a,b)=(-1)^{ab}W_N(a,b)\). In particular, \[
 Q_N(v)Q_N(w)=e^{\pi i\omega(v,w)/N}Q_N(v+w),\qquad
 \omega(v,w)=v_1w_2-v_2w_1.                                 \tag{LC.7}
\] Each \(Q_N(v)\) is unitary, \(Q_N(v)^*=Q_N(-v)\), and parity sends \(Q_N(v)\) to \(Q_N(-v)\). Set \[
 \operatorname{Op}_N(a)=\sum_{v\in\mathbb Z^2}\widehat a(v)Q_N(v),
 \quad \mathcal M_r(a)=\sum_v|\widehat a(v)|(1+|v|_1)^r.
\] The series converges absolutely in operator norm for smooth \(a\); \(\|\operatorname{Op}_N(a)\|\leq\mathcal M_0(a)\), and real symbols give Hermitian matrices. Frequency aliases cause no problem in this definition.

Since \(T\) is the identity modulo two, its finite Weyl implementer also has exact physical covariance \[
 U_NQ_N(v)U_N^*=Q_N(Tv),\qquad
 U_N\operatorname{Op}_N(a)U_N^*=\operatorname{Op}_N(a\circ M). \tag{LC.8}
\] One such implementer has the kernel \[
 U_N(x,y)=\frac{c_N}{\sqrt N}
 e_N\!\left((3x^2-2xy+3y^2)/8\right),\quad |c_N|=1,         \tag{LC.9}
\] where division by eight is modular. Orthogonality of additive characters proves unitarity; direct substitution on the two Weyl generators proves (LC.8). The kernel is parity-preserving. Its global phase can be fixed later.

\textbf{Lemma 1 (disjoint-symbol product).} If \(a,b\in C^\infty(\mathbb T^2)\) have disjoint closed supports, then for every integer \(r\geq1\), \[
 \|\operatorname{Op}_N(a)\operatorname{Op}_N(b)\|_{\rm op}
 \leq\frac{\pi^r}{r!N^r}\mathcal M_r(a)\mathcal M_r(b).       \tag{LC.10}
\]

\textbf{Proof.} Expand (LC.7) and Taylor-expand \(e^{i\pi\omega/N}\) through degree \(r-1\). For real \(z\) the integral Taylor remainder satisfies \(|e^{iz}-\sum_{j<r}(iz)^j/j!|\leq|z|^r/r!\), with no restriction on \(z\). For each \(j<r\) the coefficient sum \[
 \sum_{v,w}\widehat a(v)\widehat b(w)\omega(v,w)^j Q_N(v+w)
\] is the quantization of \[
 \frac{1}{(-4\pi^2)^j}\sum_{l=0}^j(-1)^l\binom jl
 (\partial_q^{j-l}\partial_p^l a)
 (\partial_q^l\partial_p^{j-l} b).                           \tag{LC.11}
\] All these functions vanish, including \(ab\) at \(j=0\), because derivatives are supported in the respective original closed supports. Their Fourier cancellation occurs before applying \(Q_N\), so aliases do not invalidate it. The remaining absolutely convergent sum is bounded using \(|\omega(v,w)|\leq|v|_1|w|_1\) and \(\|Q_N(v+w)\|=1\). This proves (LC.10).

For real \(a,b\), their two products are adjoints, so \[
 \|[\operatorname{Op}_N(a),\operatorname{Op}_N(b)]\|
 \leq\frac{2\pi^r}{r!N^r}\mathcal M_r(a)\mathcal M_r(b).      \tag{LC.12}
\] This all-orders statement is stronger than just Poisson commutation.

\subsection{An explicit infinite-dimensional family and a Gevrey bound}\label{locality:an-explicit-infinite-dimensional-family-and-a-gevrey-bound}

